%% ==========================================================================
%%  preamble.tex —— 《线性代数：抽象结构 · 算子及其分解》版式定义
%%  编译引擎：XeLaTeX（需要 ctex + tcolorbox + Libertinus + Noto CJK）
%% ==========================================================================

% -------------------------------------------------------------------------
% 1. 页面尺寸（B5 开本，接近国内数学教材手感）
% -------------------------------------------------------------------------
\usepackage[
  paperwidth  = 176mm,
  paperheight = 250mm,
  inner       = 21mm,
  outer       = 17mm,
  top         = 22mm,
  bottom      = 24mm,
  headheight  = 14pt,
  headsep     = 6mm,
  footskip    = 11mm,
]{geometry}

% -------------------------------------------------------------------------
% 2. 字体：西文 Libertinus，数学 Libertinus Math，中文 Noto Serif/Sans CJK SC
% -------------------------------------------------------------------------
\usepackage{fontspec}
\usepackage{unicode-math}

\setmainfont{LibertinusSerif}[
  Extension      = .otf,
  UprightFont    = *-Regular,
  ItalicFont     = *-Italic,
  BoldFont       = *-Bold,
  BoldItalicFont = *-BoldItalic,
  Numbers        = OldStyle,
  Ligatures      = TeX,
]
\setsansfont{LibertinusSans}[
  Extension      = .otf,
  UprightFont    = *-Regular,
  ItalicFont     = *-Italic,
  BoldFont       = *-Bold,
  Ligatures      = TeX,
]
\setmonofont{LibertinusMono-Regular}[Extension = .otf, Scale = MatchLowercase]
\setmathfont{LibertinusMath-Regular.otf}
% 公式里的数字用等高数字，正文用老式数字
\newfontfamily\liningfont{LibertinusSerif}[Extension=.otf, UprightFont=*-Regular,
  BoldFont=*-Bold, Numbers=Lining]

\xeCJKsetup{CJKmath = true}
\setCJKmainfont{Noto Serif CJK SC}[
  BoldFont       = Noto Serif CJK SC Bold,
  ItalicFont     = Noto Serif CJK SC,
  ItalicFeatures = {FakeSlant = 0.18},
]
\setCJKsansfont{Noto Sans CJK SC}[BoldFont = Noto Sans CJK SC Bold]
\setCJKmonofont{Noto Sans Mono CJK SC}
\newCJKfontfamily\cjkhei{Noto Sans CJK SC}[BoldFont = Noto Sans CJK SC Bold]
\newCJKfontfamily\cjklight{Noto Serif CJK SC}[BoldFont = Noto Serif CJK SC Bold]

% -------------------------------------------------------------------------
% 3. 数学、表格、图形等基础宏包
% -------------------------------------------------------------------------
\usepackage{amsmath, mathtools}   % 注意：unicode-math 已提供全部 AMS 符号，不能再载 amssymb
\usepackage{booktabs, array}
\usepackage{enumitem}
\usepackage{graphicx}
\usepackage{pifont}
\usepackage{microtype}
\usepackage{tikz}
\usetikzlibrary{calc}
\usepackage[most]{tcolorbox}
\tcbuselibrary{theorems, breakable, skins}
\usepackage{titlesec}
\usepackage{titletoc}
\usepackage{fancyhdr}
\usepackage{emptypage}

% -------------------------------------------------------------------------
% 4. 配色方案（沉静的靛蓝为主色，四种辅色区分不同类型的知识块）
% -------------------------------------------------------------------------
\usepackage{xcolor}
\definecolor{inkblue}{HTML}{16324F}      % 主色：章节标题
\definecolor{accent}{HTML}{2C5F8A}       % 辅色：节号、页眉细线
\definecolor{defnblue}{HTML}{1C6E8C}     % 定义
\definecolor{thmred}{HTML}{9B2C2C}       % 定理
\definecolor{propgreen}{HTML}{1F6F54}    % 命题
\definecolor{proofgray}{HTML}{4A5568}    % 证明思路
\definecolor{insightgold}{HTML}{8A6100}  % 要点
\definecolor{ideapurple}{HTML}{5B4B8A}   % 理解
\definecolor{pagegray}{HTML}{6B7280}

% -------------------------------------------------------------------------
% 5. 段落、行距、标点
% -------------------------------------------------------------------------
\linespread{1.30}
\setlength{\parskip}{0.28em plus 0.1em minus 0.05em}
\setlength{\parindent}{2\ccwd}
\setlength{\emergencystretch}{1.5em}
\raggedbottom
\predisplaypenalty=0        % 允许在行间公式前分页，避免大片留白
\postdisplaypenalty=0
\setlength{\abovedisplayskip}{7pt plus 2pt minus 2pt}
\setlength{\belowdisplayskip}{7pt plus 2pt minus 2pt}
\setlength{\abovedisplayshortskip}{3pt plus 2pt}
\setlength{\belowdisplayshortskip}{5pt plus 2pt}

% -------------------------------------------------------------------------
% 6. 超链接与 PDF 书签
% -------------------------------------------------------------------------
\usepackage[
  colorlinks = true,
  linkcolor  = accent,
  citecolor  = accent,
  urlcolor   = accent,
  linktoc    = all,
  bookmarksnumbered  = true,
  bookmarksopen      = true,
  bookmarksopenlevel = 1,
]{hyperref}
\usepackage{bookmark}
% PDF 元信息里含逗号与中文，必须放在 \hypersetup 里而不是宏包选项里
\hypersetup{
  pdftitle    = {线性代数：抽象结构 —— 算子及其分解},
  pdfsubject  = {线性代数讲义},
  pdfkeywords = {线性代数; 算子; 不变子空间; Jordan 标准型; 谱定理; 奇异值分解},
}

% -------------------------------------------------------------------------
% 7. 章 / 节 / 小节标题样式
% -------------------------------------------------------------------------
% 章：深色横条给出章序 + 大标题 + 双色横线
% \CTEXthechapter 与目录中的“第一章”写法保持一致
\providecommand{\CTEXthechapter}{第\,\thechapter\,章}
\titleformat{\chapter}[display]
  {\normalfont\sffamily\raggedright}
  {\colorbox{inkblue}{\parbox{\dimexpr\linewidth-2\fboxsep}{%
     \color{white}\sffamily\bfseries\normalsize
     \strut \CTEXthechapter}}}
  {6pt}
  {\huge\bfseries\color{inkblue}}
  [{\vspace{4pt}{\color{accent}\titlerule[1.1pt]}%
     \vspace{1.4pt}{\color{accent!45}\titlerule[0.5pt]}}]
\titlespacing*{\chapter}{0pt}{4pt}{22pt}

\titleformat{name=\chapter,numberless}[display]
  {\normalfont\sffamily\raggedright}
  {}
  {0pt}
  {\huge\bfseries\color{inkblue}}
  [{\vspace{4pt}{\color{accent}\titlerule[1.1pt]}%
     \vspace{1.4pt}{\color{accent!45}\titlerule[0.5pt]}}]

% 节：色条 + 编号
\titleformat{\section}[hang]
  {\normalfont\sffamily\large\bfseries\color{inkblue}}
  {\raisebox{-0.08em}{\color{accent}\rule{2.6pt}{1.05em}}\hspace{0.55em}%
   \color{accent}\thesection\hspace{0.1em}}
  {0.6em}{}
\titlespacing*{\section}{0pt}{1.7ex plus .3ex}{0.9ex}

\titleformat{\subsection}[hang]
  {\normalfont\sffamily\bfseries\color{accent!85!black}}
  {\thesubsection}{0.55em}{}
\titlespacing*{\subsection}{0pt}{1.3ex plus .2ex}{0.6ex}

% 部分（\part）：整页的大标题，编号用汉字
\renewcommand{\thepart}{\chinese{part}}
\titleformat{\part}[display]
  {\normalfont\sffamily\centering}
  {{\color{accent!45}\rule{0.22\linewidth}{0.9pt}}\\[10pt]
   \color{accent}\Large 第\,\thepart\,部分}
  {14pt}
  {\vspace{2pt}\color{inkblue}\Huge\bfseries}
  [{\vspace{12pt}{\color{accent}\rule{0.42\linewidth}{1.1pt}}}]
\titlespacing*{\part}{0pt}{80pt}{28pt}

% -------------------------------------------------------------------------
% 8. 页眉页脚
% -------------------------------------------------------------------------
\pagestyle{fancy}
\fancyhf{}
\fancyhead[LE]{\small\sffamily\color{pagegray}\leftmark}
\fancyhead[RO]{\small\sffamily\color{pagegray}\rightmark}
\fancyfoot[LE,RO]{\small\sffamily\color{accent}\liningfont\thepage}
\renewcommand{\headrulewidth}{0pt}
\renewcommand{\headrule}{{\color{accent!45}\hrule height 0.4pt}\vspace{1pt}}
\fancypagestyle{plain}{%
  \fancyhf{}%
  \renewcommand{\headrule}{}%
  \fancyfoot[LE,RO]{\small\sffamily\color{accent}\liningfont\thepage}}
\renewcommand{\chaptermark}[1]{\markboth{\CTEXthechapter\quad #1}{}}
\renewcommand{\sectionmark}[1]{\markright{\thesection\quad #1}}

% -------------------------------------------------------------------------
% 9. 目录样式
% -------------------------------------------------------------------------
\renewcommand{\contentsname}{目\hspace{1em}录}
\titlecontents{part}[0pt]
  {\addvspace{18pt}\sffamily\bfseries\color{inkblue}\large}
  {\contentspush{第\,\thecontentslabel\,部分\quad}}{}{}
\titlecontents{chapter}[0pt]
  {\addvspace{7pt}\sffamily\bfseries\color{inkblue}}
  {\contentspush{\thecontentslabel\enspace}}{}
  {\hfill\color{accent}\liningfont\contentspage}
\titlecontents{section}[2.9em]
  {\addvspace{1pt}\small}
  {\contentslabel{2.6em}}{\hspace*{-2.6em}}
  {\titlerule*[0.6pc]{\color{accent!35}.}\color{accent}\liningfont\contentspage}
\titlecontents{subsection}[5.8em]
  {\addvspace{0pt}\small\color{black!80}}
  {\contentslabel{3.1em}}{\hspace*{-3.1em}}
  {\titlerule*[0.6pc]{\color{accent!25}.}\color{accent!80}\liningfont\contentspage}

% -------------------------------------------------------------------------
% 10. 知识块：定义 / 定理 / 命题 / 证明思路 / 要点 / 理解 / 部分导读
% -------------------------------------------------------------------------
% 盒子内的段落：不缩进，靠空白分段（parbox=false 才能让段设置生效）
\newcommand{\jboxpar}{\setlength{\parindent}{0pt}\setlength{\parskip}{0.5em plus 0.1em}}

% 统一外观：隐去边框，只留左侧色条 + 极淡底色；三类共用一个“结果”计数器
\tcbset{
  jstatement/.style = {
    enhanced jigsaw, breakable, sharp corners, boxrule = 0pt, frame hidden,
    parbox = false,
    left = 10pt, right = 10pt, top = 7pt, bottom = 7pt,
    before skip = 11pt plus 2pt, after skip = 11pt plus 2pt,
    fonttitle = \sffamily\bfseries,
    theorem style = standard,
    separator sign none,
    terminator sign none,
    description delimiters = {\kern-0.42em（}{）},
    before upper = {\jboxpar},
  },
}

\newtcbtheorem[number within = chapter, list inside = jdeflist]{jdefn}{定义}{
  jstatement,
  colback = defnblue!5!white,
  coltitle = defnblue!88!black,
  borderline west = {2.6pt}{0pt}{defnblue},
}{def}

\newtcbtheorem[use counter from = jdefn, list inside = jthmlist]{jthm}{定理}{
  jstatement,
  colback = thmred!5!white,
  coltitle = thmred!88!black,
  borderline west = {2.6pt}{0pt}{thmred},
}{thm}

\newtcbtheorem[use counter from = jdefn, list inside = jproplist]{jprop}{命题}{
  jstatement,
  colback = propgreen!5!white,
  coltitle = propgreen!88!black,
  borderline west = {2.6pt}{0pt}{propgreen},
}{prop}

% 证明思路：无底色，细灰竖线，末尾一个小方块（用 \rule 画，不依赖字体字形）
\newcommand{\jqed}{\textcolor{proofgray!80}{\rule[0.03em]{0.46em}{0.46em}}}
\newtcolorbox{jproof}[1]{
  enhanced jigsaw, breakable, sharp corners, boxrule = 0pt, frame hidden,
  parbox = false,
  colback = white,
  borderline west = {1pt}{0pt}{proofgray!45},
  left = 9pt, right = 6pt, top = 4pt, bottom = 4pt,
  before skip = 9pt plus 2pt, after skip = 11pt plus 2pt,
  fontupper = \small,
  before upper = {\jboxpar\textcolor{proofgray}{\sffamily\bfseries #1\,}\hspace{0.45em}},
  after upper  = {\unskip\nobreak\hfill\jqed},
}

% 要点（原 Markdown 的引用块）
\newtcolorbox{jinsight}{
  enhanced jigsaw, breakable, sharp corners, boxrule = 0pt, frame hidden,
  parbox = false,
  colback = insightgold!8!white,
  borderline west = {2.6pt}{0pt}{insightgold!85},
  left = 10pt, right = 10pt, top = 7pt, bottom = 7pt,
  before skip = 11pt plus 2pt, after skip = 11pt plus 2pt,
  before upper = {\jboxpar\textcolor{insightgold!85!black}{\ding{72}\ \sffamily\bfseries 要点\,}\hspace{0.4em}},
}

% 理解 / 构造目标一类的旁注
\newtcolorbox{jidea}[1]{
  enhanced jigsaw, breakable, sharp corners, boxrule = 0pt, frame hidden,
  parbox = false,
  colback = ideapurple!5!white,
  borderline west = {2.6pt}{0pt}{ideapurple!85},
  left = 10pt, right = 10pt, top = 6pt, bottom = 6pt,
  before skip = 10pt plus 2pt, after skip = 10pt plus 2pt,
  fontupper = \small,
  before upper = {\jboxpar\textcolor{ideapurple!85!black}{\sffamily\bfseries #1\,}\hspace{0.45em}},
}

% 每一“部分”的导读：整页宽的浅色大盒子
\newtcolorbox{partintro}[1]{
  enhanced, breakable, sharp corners,
  parbox = false,
  colback = accent!4!white, colframe = accent!35,
  boxrule = 0.6pt,
  borderline west = {3pt}{0pt}{accent},
  left = 14pt, right = 14pt, top = 14pt, bottom = 12pt,
  before skip = 6pt, after skip = 16pt,
  fonttitle = \sffamily\bfseries,
  coltitle = white,
  colbacktitle = accent,
  title = {\ #1},
  before upper = {\jboxpar},
  attach boxed title to top left = {xshift = 6pt, yshift = -\tcboxedtitleheight/2},
  boxed title style = {sharp corners, boxrule = 0pt, colback = accent},
}

% -------------------------------------------------------------------------
% 11. 列表样式
% -------------------------------------------------------------------------
\newcommand{\jsquare}{\textcolor{accent}{\raisebox{0.22em}{\rule{0.42em}{0.42em}}}}
\newlist{jitem}{itemize}{3}
\setlist[jitem]{label = \jsquare,
  leftmargin = 1.9em, itemsep = 1pt, topsep = 4pt, parsep = 2pt}
\newlist{jenum}{enumerate}{3}
\setlist[jenum]{label = {\textcolor{accent}{\sffamily\bfseries\arabic*.}},
  leftmargin = 2.1em, itemsep = 1pt, topsep = 4pt, parsep = 2pt}

% -------------------------------------------------------------------------
% 12. 常用数学记号
% -------------------------------------------------------------------------
\newcommand{\FF}{\mathbb{F}}
\newcommand{\RR}{\mathbb{R}}
\newcommand{\CC}{\mathbb{C}}
\DeclareMathOperator{\spann}{span}
\DeclareMathOperator{\diag}{diag}

% 扉页上用的装饰线
\newcommand{\titlerulebig}[1]{{\color{#1}\rule{\linewidth}{1.2pt}}}
